% Nouveau EGA V (cote 25): archive cote 25, site page 27 (PDF page 28).
% typescript. Draft compared: none.
% AI-transcribed by Claude Opus 5.5 (claude-opus-5-5, adaptive thinking, effort high), Message Batches API, batch ega25b (msgbatch_01A7n5YuXubUPNK1zsepi8RG), 2026-11-13; original AI transcription, subsequently checked against both supplied scan parts; see PAGES_CHECKED.tsv and CHECK_DECISIONS.tsv.
\scanpage{25/27}
\mspage{\hand{23}}

\var{\deleted{Remords au N°4 . Les résultats 4.6. b) et c) sont enti\var{\xout{\unsure{r|m}}}{}{typewriter strike-out}{ega25_25_027-v0002}èrement
contenus dans 4.7. b) et c) ; j'avais cru une telle redite justifiable du
point de vue démonstration, celle de 4.7. me semblant mettre en oeuvre des
résultats nettement plus délicats (sous forme de 15.2.2.). Or, tant qu'on
ne tient qu'à prouver dans 4.7. b) de c) que U est ouvert, (à l'exclusion
du fait $\underline{G}$ U plat \var{\xout{xxx}}{}{typewriter strike-out}{ega25_25_027-v0003} rel à $\overset{\hand{\vee}}{P}$ ), on n'a nul besoin de 15.2.2. , mais la
méthode habituelle (utilisée au par.12) par le critère valuatif donne le
résultat voulu de façon bien simple. Je suggère donc:a\var{\xout{\unsure{x}}}{}{typewriter strike-out}{ega25_25_027-v0004}) de vider les
parties b) et c) de 4.6. ; b) d'énoncer 4.7. b) et c) avec la seule
conclusion U resp. E ouvert, en utilisant le critère valuatif (bien entendu,
4.7. a) est vidé de l'énoncé), c) d'ajouter une autre proposition\var{\inserted[typed]{4.8.}}{4.8.}{insertion}{ega25_25_027-v0005}, ou un
corollaire, disant que sous les conditions de 4.7. $\underline{G}$ U est plat rel à $\overset{\hand{\vee}}{P}$ ,
ce qui se démontre à l'aide de 15.2.2.}}{}{deletion}{ega25_25_027-v0001}

\var{\over{\unsure{5|6}}{7}}{7}{writing over}{ega25_25_027-v0006}. \uline{Application à la construct\var{\xout{i}}{}{typewriter strike-out}{ega25_25_027-v0007}ion de sections hyperplanes ou de multisections
de type précisé}.

\var{\over{\unsure{5|6}}{7}}{7}{writing over}{ega25_25_027-v0008}.1. Notons que si S=Spec(k), ou k est un corps \uline{infini}, alors tout ouvert
non vide de $\overset{\hand{\vee}}{P}$ contient un point rationnel sur k; donc avec les notations de
\var{\xout{\unsure{xxx}}}{}{typewriter strike-out}{ega25_25_027-v0009}4.1. , si E (défini en termes d'une propriété constructible P) contient
le point générique \hand{$\eta$} , il cotient \var{\xout{au moins}}{}{typewriter strike-out}{ega25_25_027-v0010} un point rationnel sur k, donc
il existe une section hyperplane de X (définie sur k) ayant la propriété P.
D'autre part, S étant de nouveau quelconque, il est immédiat que pour tout
s\hand{$\in$}S et tout point\var{\inserted{$\xi_0$}}{ $\xi_0$}{insertion}{ega25_25_027-v0011} de la fibre $\overset{\hand{\vee}}{P}_s$ rationnel sur k(s), il existe une section
\hand{$\xi$} de $\overset{\hand{\vee}}{P}$ sur un voisinage ouvert U de s qui passe par \hand{$\xi$}$_0$ . Si alors E est
encore \var{\xout{xxx}}{}{typewriter strike-out}{ega25_25_027-v0012} défini\var{\xout{x}}{}{typewriter strike-out}{ega25_25_027-v0013} comme dans 4.1. en termes d'une propriété constructible
P, et si on dispose (p ex grâce \var{\deleted{à 5.\unsure{5}.}}{}{deletion}{ega25_25_027-v0014}\var{\inserted{au N°5}}{au N°5}{insertion}{ega25_25_027-v0015}) du fait que E est ouvert, alors
si \hand{$\xi$}$_0$\hand{$\in$} E, \var{\xout{on peut prendre M}}{}{typewriter strike-out}{ega25_25_027-v0016} la section \hand{$\xi$} est une section de E sur \var{\xout{\unsure{Sx}}}{}{typewriter strike-out}{ega25_25_027-v0017}
U, du moins si on rapetisse suffisamment U. Donc, on aura construit une
section hyperplane\var{\inserted[typed]{$Y_{\hand{\xi}}$}}{ $Y_{\hand{\xi}}$}{insertion}{ega25_25_027-v0018} de X , au dessus d'un voisinage ouvert U de s, telle que

\var{\margin{27}}{\placedmargin{27}}{marginal note placed}{ega25_25_027-v0019}
